\documentclass[10pt]{article}
\usepackage{titlesec}
\usepackage{geometry}
\geometry{verbose,tmargin=.9in,bmargin=.9in,lmargin=1.0in,rmargin=1.0in}
\usepackage{amsmath,amsfonts,amsthm,amssymb}
\usepackage{url}
\usepackage{color}
\usepackage[usenames,dvipsnames,svgnames,table]{xcolor}
\usepackage[colorlinks=true, linkcolor=red, urlcolor=blue, citecolor=gray]{hyperref}
\usepackage{float}
\usepackage{caption}
\usepackage{subcaption}
\usepackage{graphicx}
\usepackage{wrapfig}
\usepackage{booktabs}
\usepackage{longtable}
\usepackage{enumitem}
\usepackage{multicol}
\usepackage{etoolbox}
%\usepackage{bbm}
\usepackage{ulem}

\newcommand{\bs}[1]{\boldsymbol{#1}}
\newcommand{\bv}[1]{\mathbf{#1}}
\newcommand{\R}{\mathbb{R}}
\newcommand{\E}{\mathbb{E}}
\DeclareMathOperator*{\argmin}{arg\,min}
\DeclareMathOperator*{\argmax}{arg\,max}
\DeclareMathOperator{\rank}{rank}
\DeclareMathOperator{\Var}{Var}
\DeclareMathOperator{\Cost}{Cost}
\DeclareMathOperator{\cut}{cut}
\DeclareMathOperator{\tr}{tr}

\definecolor{nyuDarkPurple}{HTML}{330662}
\definecolor{nyuOfficialPurple}{HTML}{57068c}

\newcommand{\spara}[1]{\vspace{.5em}\noindent {\large\sffamily\textcolor{nyuOfficialPurple}{#1}}}
\titleformat{\section}[hang]{\Large\sffamily\color{nyuDarkPurple}}{\thesection}{1em}{}
\titleformat{\subsection}[hang]{\large\sffamily\color{nyuDarkPurple}}{\thesection}{1em}{}
\titleformat{\subsubsection}[hang]{\normalsize\sffamily\color{gray}}{\thesection}{1em}{}

\usepackage{fancyhdr}
\setlength{\headheight}{30pt}
\pagestyle{fancy}
\lhead{}
\rhead{\thepage}
\pagenumbering{gobble}

\setcounter{secnumdepth}{0}

\begin{document}
	
\begin{center}
	\normalsize
	New York University Tandon School of Engineering
	
	Computer Science and Engineering
	\medskip
	
	\large
	CS-GY 6763: Homework 3. 
	
	\medskip
	
	\normalsize 
	\noindent \emph{Collaboration is allowed on this problem set, but solutions must be written-up individually. Please list collaborators for each problem separately, or write ``No Collaborators" if you worked alone.}
	\medskip
\end{center} 



\subsection{Problem 1: Compressed classification.}
\textbf{(10 pts)} In machine learning, the goal of many classification methods (like support vector machines) is to separate data into classes using a \emph{separating hyperplane}.

Recall that a hyperplane in $\R^d$ is defined by a unit vector $a \in \R^d$ ($\|a\|_2 = 1$) and scalar $c \in \R$. It contains all  $h \in \R^d$ such that $\langle a, h\rangle = c$. 

Suppose our dataset consists of $n$ unit vectors in $\R^d$ (i.e., each data point is normalized to have norm $1$). These points can be separated into two sets $X, Y$,
with the guarantee that there exists a hyperplane such that every point in $X$ is on one side and every point in 
$Y$ is on the other. In other words, for all $x\in X, \langle a, x\rangle > c$ and for all $y\in Y, \langle a, y\rangle < c$.

Furthermore, suppose that the $\ell_2$ distance of each point in $X$ and $Y$ to this separating hyperplane is at least $\epsilon$. When this is the case, the hyperplane is said to have ``margin'' $\epsilon$. 

\begin{enumerate}
	\item Show that this margin assumption equivalently implies that for all $x\in X, \langle a, x\rangle \geq c + \epsilon$ and for all $y\in Y, \langle a, y\rangle \leq c - \epsilon$.
	
	\item Show that if we use a Johnson-Lindenstrauss map $\Pi$ to reduce our data points to $O(\log n/\epsilon^2)$ dimensions, then the dimension reduced data can still be separated by a hyperplane with margin $\epsilon/4$, with high probability (say $> 9/10$).
\end{enumerate}
 



\subsection{Problem 2: Concentration of Random Vectors}
\textbf{(12 pts)} In Stochastic Gradient Descent, we replace the true gradient vector with a stochastic gradient that is equal to the true gradient in expectation. Our analysis in class \emph{only used} equality in expectation, although more refined analysis of SGD often requires understanding how well the stochastic gradient \emph{concentrates} around its expectation. Previously, all concentration results we studied apply to random numbers.  For this problem, you will prove a basic concentration inequality for random vectors. 

In particular, let $\bv{x}_1, \ldots, \bv{x}_k \in \R^d$ be i.i.d. random vectors in $d$ dimensions (independent, drawn from the same distribution) with mean $\bs{\mu}$. I.e., $\E[\bv{x}_i] = \bs{\mu}$. Further suppose that $\E\left[\|\bv{x}_i - \bs{\mu}\|_2^2\right] = \sigma^2$. $\sigma^2$ is a natural generalization of ``variance'' to a random vector. Let $\bv{s} = \frac{1}{k}\sum_{i=1}^k\bv{x}_i$. Prove that if $k \geq O\left(\frac{1/\delta}{\epsilon^2}\right)$ then
\begin{align*}
	\Pr\left[\|\bv{s} - \bs{\mu}\right\|_2 \geq \epsilon \sigma] \leq \delta.
\end{align*}

% --- PS3 Problem 2: Concentration of random vectors (clean writeup) ---

\subsection*{Problem 2: Concentration of random vectors}

\textbf{Given.}
Let $x_1,\dots,x_k \in \mathbb{R}^d$ be i.i.d.\ random vectors such that
\[
\E[x_i]=\mu,
\qquad
\E\!\left[\|x_i-\mu\|_2^2\right]=\sigma^2.
\]
Define the sample mean
\[
s=\frac{1}{k}\sum_{i=1}^k x_i.
\]

\textbf{Sanity check (as $k\to\infty$).}
By linearity of expectation,
\[
\E[s]=\E\!\left[\frac{1}{k}\sum_{i=1}^k x_i\right]
=\frac{1}{k}\sum_{i=1}^k \E[x_i]
=\mu.
\]
So the average is centered at $\mu$; we now quantify how large $k$ must be so that $s$ concentrates around $\mu$.

\textbf{Goal.}
Show that
\[
\Pr\!\left(\|s-\mu\|_2 \ge \varepsilon\sigma\right)\le \delta
\qquad
\text{as long as }
k = O\!\left(\frac{1}{\delta\varepsilon^2}\right).
\]

\medskip
\textbf{Step 1: Center the variables.}
Let
\[
v_i := x_i-\mu.
\]
Then
\[
\E[v_i]=\E[x_i-\mu]=\mu-\mu=0,
\qquad
\E\!\left[\|v_i\|_2^2\right]=\E\!\left[\|x_i-\mu\|_2^2\right]=\sigma^2.
\]
Also,
\[
s-\mu
=
\frac{1}{k}\sum_{i=1}^k (x_i-\mu)
=
\frac{1}{k}\sum_{i=1}^k v_i.
\]

\medskip
\textbf{Step 2: Compute $\E\|s-\mu\|_2^2$.}
Start with
\[
\E\|s-\mu\|_2^2
=
\E\left\|\frac{1}{k}\sum_{i=1}^k v_i\right\|_2^2
=
\frac{1}{k^2}\E\left\|\sum_{i=1}^k v_i\right\|_2^2.
\]
Expand the squared norm:
\begin{align*}
\left\|\sum_{i=1}^k v_i\right\|_2^2
&=
\left\langle \sum_{i=1}^k v_i,\ \sum_{j=1}^k v_j \right\rangle \\
&=
\sum_{i=1}^k \|v_i\|_2^2
+
2\sum_{i<j}\langle v_i,v_j\rangle.
\end{align*}
Taking expectation and using linearity,
\[
\E\left\|\sum_{i=1}^k v_i\right\|_2^2
=
\sum_{i=1}^k \E\|v_i\|_2^2
+
2\sum_{i<j}\E\langle v_i,v_j\rangle.
\]
For $i\neq j$, $v_i$ and $v_j$ are independent and mean-zero, so
\[
\E\langle v_i,v_j\rangle
=
\left\langle \E[v_i],\E[v_j]\right\rangle
=
\langle 0,0\rangle
=0.
\]
Thus the cross terms vanish and
\[
\E\left\|\sum_{i=1}^k v_i\right\|_2^2
=
\sum_{i=1}^k \sigma^2
=
k\sigma^2.
\]
Therefore,
\[
\boxed{
\E\|s-\mu\|_2^2
=
\frac{1}{k^2}\cdot k\sigma^2
=
\frac{\sigma^2}{k}.
}
\]

\medskip
\textbf{Step 3: Apply Markov to $\|s-\mu\|_2^2$.}
Let
\[
Y := \|s-\mu\|_2^2 \quad(\ge 0).
\]
Then
\[
\Pr\!\left(\|s-\mu\|_2 \ge \varepsilon\sigma\right)
=
\Pr\!\left(\|s-\mu\|_2^2 \ge \varepsilon^2\sigma^2\right)
=
\Pr\!\left(Y \ge \varepsilon^2\sigma^2\right).
\]
By Markov's inequality,
\[
\Pr(Y\ge a)\le \frac{\E[Y]}{a}.
\]
Set $a=\varepsilon^2\sigma^2$ and use $\E[Y]=\sigma^2/k$:
\[
\Pr\!\left(Y \ge \varepsilon^2\sigma^2\right)
\le
\frac{\E[Y]}{\varepsilon^2\sigma^2}
=
\frac{\sigma^2/k}{\varepsilon^2\sigma^2}
=
\frac{1}{k\varepsilon^2}.
\]
So
\[
\boxed{
\Pr\!\left(\|s-\mu\|_2 \ge \varepsilon\sigma\right)
\le
\frac{1}{k\varepsilon^2}.
}
\]

\medskip
\textbf{Step 4: Choose $k$ so the failure probability is $\le \delta$.}
We want
\[
\frac{1}{k\varepsilon^2}\le \delta
\quad\Longrightarrow\quad
k \ge \frac{1}{\delta\varepsilon^2}.
\]
Hence it suffices to take
\[
\boxed{
k = O\!\left(\frac{1}{\delta\varepsilon^2}\right),
}
\]
which proves the claim. $\square$



\subsection{Problem 3: Gradient Descent with Decaying Step-size}
\textbf{(10 pts)}
In class we showed that gradient descent with step size $\eta = R/G\sqrt{T}$ converges to an $\epsilon$ approximate minimizer of a convex $G$-Lipschitz function in $T = R^2G^2/\epsilon^2$ steps if our starting point $\bv{x}^{(0)}$ satisfies $\|\bv{x}^{(0)} - \bv{x}^*\|_2 \leq R$. 
Choosing this step size requires knowing $G$, $R$ and moreover $T$ in advance, which might not be reasonable in a lot of settings. For example, when training machine learning models, we might not be able to estimate how long it will take to reach a point where test accuracy levels off. Instead, we want to be able to keep running the algorithm, achieving better and better accuracy as we do. 

Here, we analyze a variant of gradient descent with a variable step size that avoids this limitation. In particular, consider running gradient descent with the update $\bv{x}^{(i+1)} = \bv{x}^{(i)}- \eta \nabla f(\bv{x}^{(i)})$, where
\begin{align*}
	\eta = \frac{f(\bv{x}^{(i)}) - f(\bv{x}^{*})}{\|\nabla f(\bv{x}^{(i)})\|_2^2}.
\end{align*}
This step size requires knowledge of $f(\bv{x}^{*})$, but not $\bv{x}^*$, which may be reasonable in some settings. Moreover, since it's just one parameter, grid search can be more easily used to ``guess'' $f(\bv{x}^{*})$ than the three parameters $G,R,T$. More complex approaches can remove the need to know this value entirely.
Prove that, if we run gradient descent for $T = O(R^2G^2/\epsilon^2)$ steps using the step size above then $\hat{\bv{x}} = \min_{i \in 0, \ldots, T}f(\bv{x}^{(i)})$ satisfies $f(\hat{\bv{x}}) \leq f(\bv{x}^*) + \epsilon$. 

\textbf{Hint:} Prove that our distance from the optimum $\|\bv{x}^{(i)} - \bv{x}^*\|_2$ {always decreases} with this choice of step size, and the decrease is larger if our gap from the objective value $f(\bv{x}^{(i)}) - f(\bv{x}^{*})$ is larger.


\subsection{Problem 4: Separation Oracles}
\textbf{(12 pts)} Describe efficient separation oracles for each of the following families of convex sets.  Here, ``efficient'' means linear time plus $O(1)$ calls to any additional oracles provided to you.

\begin{enumerate}[label=(\alph*)]
	\item
	The set $A \cap B$, given separation oracles for $A$ and $B$.
	
	\item
	The $\ell_1$ ball: $\| \bv{x} \|_1 \leq 1$.
	
	\item
	Any convex set $A$, given a \emph{projection oracle} for $A$. Recall that a projection oracle, given a point $\bv{x}$, returns
	\begin{equation*}
		\text{Proj}_A(\bv{x}) =  \argmin_{y \in A} \| \bv{x} - \bv{y} \|_2.
	\end{equation*}
	
\end{enumerate}
Above you may wish to use the following fact that was stated but not formally proven in class: for any point $\bv{x}$, convex set $A$, and $\bv{z}\in A$, $\|\bv{z}-\text{Proj}_A(\bv{x})\|_2 \leq \|\bv{z} - \bv{x}\|_2$.







\end{document}